% Figure: round-trip energy accounting for the CAES case study, two columns.
% Left bar: diabatic plant (charge work + fuel reheat in, electricity out).
% Right bar: adiabatic plant (charge work in, stored heat returned, no fuel).
% Heights drawn to scale against the charge work taken as 100 units.
\begin{figure}[htbp]
\centering
\begin{tikzpicture}[font=\small, scale=1.0]
  % ---- diabatic plant (left) ----
  % charge work in: 100 units, height 3.0, x 0..1.6
  \fill[englime!35] (0,0) rectangle (1.6,3.0);
  \node[font=\scriptsize, anchor=south, align=center] at (0.8,3.05) {charge\\work $100$};
  % fuel reheat added on discharge: 120 units stacked above, height 3.6
  \fill[engamber!40] (0,3.0) rectangle (1.6,6.6);
  \node[font=\scriptsize, anchor=south, align=center] at (0.8,6.65) {fuel reheat\\$120$};
  % outputs to the right of the diabatic bar
  % electricity out: 95 units, height 2.85
  \fill[engblue!35] (2.2,0) rectangle (3.4,2.85);
  \node[engblue, font=\scriptsize, anchor=west, align=left] at (3.45,1.4) {electricity\\out $95$};
  % losses: 125 units, height 3.75
  \fill[enggray!35] (2.2,2.85) rectangle (3.4,6.6);
  \node[enggray, font=\scriptsize, anchor=west, align=left] at (3.45,4.7)
    {stack heat,\\charge-heat dumped,\\irreversibility $125$};
  \node[engblack, font=\scriptsize, anchor=north, align=center] at (0.8,-0.3)
    {diabatic\\$\eta_{\text{rt}} \approx 0.43$};
  % ---- adiabatic plant (right) ----
  \begin{scope}[xshift=7.2cm]
    % charge work in: 100 units, height 3.0
    \fill[englime!35] (0,0) rectangle (1.6,3.0);
    \node[font=\scriptsize, anchor=south, align=center] at (0.8,3.05) {charge\\work $100$};
    % stored compression heat returned on discharge: drawn as recycled, no new primary energy
    \fill[engred!25] (0,3.0) rectangle (1.6,4.2);
    \node[font=\scriptsize, anchor=south, align=center] at (0.8,4.25) {stored heat\\returned};
    % outputs: electricity 70 units height 2.1, losses 30 units height 0.9
    \fill[engblue!35] (2.2,0) rectangle (3.4,2.1);
    \node[engblue, font=\scriptsize, anchor=west, align=left] at (3.45,1.05) {electricity\\out $70$};
    \fill[enggray!35] (2.2,2.1) rectangle (3.4,3.0);
    \node[enggray, font=\scriptsize, anchor=west, align=left] at (3.45,2.55) {losses $30$};
    \node[engblack, font=\scriptsize, anchor=north, align=center] at (0.8,-0.3)
      {adiabatic\\$\eta_{\text{rt}} \approx 0.70$};
  \end{scope}
\end{tikzpicture}
\caption[Round-trip energy accounting for two CAES plant types]{Round-trip
energy accounting for the two compressed-air storage architectures of the
case study, each bar normalised to one hundred units of charge work. The
diabatic plant (left) returns more electricity than its charge work because
it burns fuel to reheat the air before expansion; counted against the total
energy in, charge work plus fuel, its round-trip figure is the modest value
quoted at the foot of the bar, and a large band leaves as dumped
compression heat, stack heat, and irreversibility. The adiabatic plant
(right) burns no fuel: it returns the stored compression heat to the air
before expansion, so its only primary input is the charge work, and a
larger fraction of that work comes back as electricity. The comparison is
the first law deciding the denominator: counting fuel as input penalises
the diabatic plant, and storing the heat instead of buying it back is what
lifts the adiabatic round-trip efficiency.}
\label{fig:vol04ch01:caes-sankey}
\end{figure}
